\documentclass{beamer}
\usetheme{AnnArbor}

\usepackage[T1]{fontenc}
\usepackage[utf8]{inputenc}
\usepackage[spanish]{babel}
\usepackage{amsmath}
\usepackage{tikz}
\usetikzlibrary{arrows.meta,calc}

\title{Fundamentos Matemáticos de la Animación}
\subtitle{Geometría, tiempo y movimiento en 2D}
\author{HCI:Animación}
\date{Mayo 2026}

\begin{document}

%------------------------------------------------
\begin{frame}
\titlepage
\end{frame}

%------------------------------------------------
\begin{frame}{Idea central del curso}
\begin{block}{Mensaje fundamental}
La animación no es simplemente dibujar figuras.
\end{block}

\medskip

La animación consiste en hacer que una o varias propiedades cambien con el tiempo.

\[
\text{Animación} = \text{Geometría} + \text{Tiempo}
\]

\medskip

Por eso, animar implica trabajar con:

\begin{itemize}
    \item Funciones del tiempo.
    \item Vectores.
    \item Movimiento.
    \item Integración numérica.
    \item Oscilaciones.
    \item Rotaciones.
\end{itemize}
\end{frame}

%------------------------------------------------
\begin{frame}{De dibujo estático a animación}
Un dibujo estático tiene valores fijos:

\[
x = 300,\qquad y = 200
\]

\medskip

Una animación tiene valores que dependen del tiempo:

\[
x = x(t),\qquad y = y(t)
\]

\medskip

\begin{block}{Idea clave}
Un objeto se anima cuando alguna propiedad se convierte en una función del tiempo.
\end{block}
\end{frame}

%------------------------------------------------
\begin{frame}{Ejemplo visual: cambio en el tiempo}
\begin{center}
\begin{tikzpicture}[scale=0.85]
    \draw[->] (0,0) -- (7,0) node[right] {Tiempo};
    \draw[->] (0,0) -- (0,3) node[above] {Posición};

    \filldraw[blue] (1,0.8) circle (3pt);
    \filldraw[blue] (3,1.5) circle (3pt);
    \filldraw[blue] (5,2.2) circle (3pt);

    \draw[thick,blue] (1,0.8) -- (3,1.5) -- (5,2.2);

    \node[below] at (1,0) {$t_0$};
    \node[below] at (3,0) {$t_1$};
    \node[below] at (5,0) {$t_2$};

    \node[right] at (5,2.2) {$x(t)$};
\end{tikzpicture}
\end{center}

\medskip

La posición no se define una sola vez.  
Se actualiza conforme avanza el tiempo.
\end{frame}

%------------------------------------------------
\begin{frame}{Agenda}
Analizaremos los siguientes fundamentos matemáticos:

\begin{enumerate}
    \item Funciones del tiempo.
    \item Movimiento uniforme.
    \item Integración numérica.
    \item Movimiento oscilatorio.
    \item Movimiento circular.
    \item Rotación en 2D.
    \item Vectores.
\end{enumerate}

\medskip

Cada tema conecta directamente con animación en computadora.
\end{frame}

%================================================
\section{Funciones del tiempo}

%------------------------------------------------
\begin{frame}{1. Funciones del tiempo}
En animación, una propiedad se modela como:

\[
p(t)
\]

donde:

\begin{itemize}
    \item \(p\) es una propiedad del objeto.
    \item \(t\) es el tiempo.
    \item \(p(t)\) indica el valor de esa propiedad en el instante \(t\).
\end{itemize}

\medskip

Ejemplos:

\[
x(t),\quad y(t),\quad \theta(t),\quad r(t)
\]
\end{frame}

%------------------------------------------------
\begin{frame}{Propiedades animables}
Una función del tiempo puede controlar:

\begin{itemize}
    \item Posición: \(x(t), y(t)\).
    \item Rotación: \(\theta(t)\).
    \item Escala: \(s(t)\).
    \item Radio: \(r(t)\).
    \item Color: \(C(t)\).
\end{itemize}

\medskip

\begin{block}{Interpretación}
Animar significa decidir cómo cambia una propiedad conforme avanza \(t\).
\end{block}
\end{frame}

%------------------------------------------------
\begin{frame}{Ejemplo matemático simple}
Supongamos que la posición horizontal es:

\[
x(t)=100+50t
\]

\medskip

Entonces:

\begin{itemize}
    \item En \(t=0\): \(x(0)=100\).
    \item En \(t=1\): \(x(1)=150\).
    \item En \(t=2\): \(x(2)=200\).
\end{itemize}

\medskip

El objeto se mueve porque \(x\) cambia con el tiempo.
\end{frame}

%================================================
\section{Movimiento uniforme}

%------------------------------------------------
\begin{frame}{2. Movimiento uniforme}
El movimiento uniforme ocurre cuando la velocidad es constante.

\[
x(t)=x_0+vt
\]

donde:

\begin{itemize}
    \item \(x_0\) es la posición inicial.
    \item \(v\) es la velocidad.
    \item \(t\) es el tiempo.
    \item \(x(t)\) es la posición en el tiempo \(t\).
\end{itemize}
\end{frame}

%------------------------------------------------
\begin{frame}{Interpretación de la velocidad}
La velocidad mide el cambio de posición por unidad de tiempo.

\[
v=\frac{\Delta x}{\Delta t}
\]

\medskip

Si:

\[
v=200\; \text{píxeles/segundo}
\]

entonces el objeto avanza 200 píxeles en un segundo.

\medskip

\begin{block}{Idea clave}
La velocidad no es solo un número; tiene unidades.
\end{block}
\end{frame}

%------------------------------------------------
\begin{frame}{Movimiento uniforme en código}
En una animación por frames no usamos directamente:

\[
x(t)=x_0+vt
\]

\medskip

Usamos una actualización cuadro a cuadro:

\[
x_{\text{nuevo}}=x_{\text{actual}}+v\cdot dt
\]

donde:

\begin{itemize}
    \item \(dt\) es el tiempo entre frames.
    \item \(v\cdot dt\) es el avance en ese frame.
\end{itemize}
\end{frame}

%------------------------------------------------
\begin{frame}{Ejemplo numérico}
Supongamos:

\[
v=200\; \text{píxeles/segundo}
\]

\[
dt=0.016\; \text{segundos}
\]

Entonces:

\[
\Delta x=v\cdot dt
\]

\[
\Delta x=200(0.016)=3.2
\]

\medskip

El objeto avanza \(3.2\) píxeles en ese frame.
\end{frame}

%================================================
\section{Integración numérica}

%------------------------------------------------
\begin{frame}{3. Integración numérica}
El modelo continuo del movimiento uniforme es:

\[
\frac{dx}{dt}=v
\]

\medskip

Pero el computador trabaja en pasos discretos:

\[
t_0,\; t_1,\; t_2,\; \ldots
\]

\medskip

Por eso aproximamos la evolución con:

\[
x_{n+1}=x_n+v\Delta t
\]
\end{frame}

%------------------------------------------------
\begin{frame}{De derivada a fórmula discreta}
Partimos de la aproximación:

\[
\frac{dx}{dt}
\approx
\frac{x_{n+1}-x_n}{\Delta t}
\]

Como:

\[
\frac{dx}{dt}=v
\]

entonces:

\[
\frac{x_{n+1}-x_n}{\Delta t}
\approx v
\]

Despejando:

\[
x_{n+1}=x_n+v\Delta t
\]
\end{frame}

%------------------------------------------------
\begin{frame}{Interpretación acumulativa}
Cada frame suma un pequeño avance:

\[
\Delta x_0=v\Delta t_0
\]

\[
\Delta x_1=v\Delta t_1
\]

\[
\Delta x_2=v\Delta t_2
\]

\medskip

Después de varios frames:

\[
x_n=x_0+\sum_{k=0}^{n-1}v\Delta t_k
\]
\end{frame}

%------------------------------------------------
\begin{frame}{Relación con la integral}
La suma discreta aproxima una integral:

\[
x(t)=x_0+\int_0^t v(\tau)\,d\tau
\]

\medskip

Si \(v\) es constante:

\[
x(t)=x_0+v\int_0^t d\tau
\]

\[
x(t)=x_0+vt
\]

\medskip

El código aproxima esta solución cuadro a cuadro.
\end{frame}

%================================================
\section{Movimiento oscilatorio}

%------------------------------------------------
\begin{frame}{4. Movimiento oscilatorio}
Un movimiento oscilatorio va y viene alrededor de una posición central.

\[
m(t)=A\sin(\omega t)
\]

donde:

\begin{itemize}
    \item \(m(t)\) es el desplazamiento oscilatorio.
    \item \(A\) es la amplitud.
    \item \(\omega\) es la velocidad angular.
    \item \(t\) es el tiempo.
\end{itemize}
\end{frame}

%------------------------------------------------
\begin{frame}{Amplitud}
La función seno cumple:

\[
-1 \leq \sin(\omega t) \leq 1
\]

Por tanto:

\[
-A \leq A\sin(\omega t) \leq A
\]

\medskip

\begin{block}{Interpretación}
La amplitud \(A\) define el máximo desplazamiento respecto al centro.
\end{block}
\end{frame}

%------------------------------------------------
\begin{frame}{Movimiento oscilatorio en animación}
La misma ecuación puede controlar muchas propiedades:

\[
m(t)=A\sin(\omega t)
\]

\medskip

Ejemplos:

\[
x(t)=x_0+A\sin(\omega t)
\]

\[
y(t)=y_0+A\sin(\omega t)
\]

\[
r(t)=r_0+A\sin(\omega t)
\]

\medskip

La función es la misma; cambia la propiedad que controlamos.
\end{frame}

%------------------------------------------------
\begin{frame}{Gráfica conceptual del seno}
\begin{center}
\begin{tikzpicture}[scale=0.85]
    \draw[->] (0,0) -- (7,0) node[right] {$t$};
    \draw[->] (0,-1.6) -- (0,1.6) node[above] {$m(t)$};

    \draw[thick,blue,domain=0:6.28,samples=100]
        plot (\x,{sin(\x r)});

    \draw[dashed] (0,1) -- (6.7,1);
    \draw[dashed] (0,-1) -- (6.7,-1);

    \node[left] at (0,1) {$A$};
    \node[left] at (0,-1) {$-A$};
    \node[below] at (3.14,0) {$\pi$};
\end{tikzpicture}
\end{center}

\[
m(t)=A\sin(\omega t)
\]

El movimiento es suave, periódico y acotado.
\end{frame}

%================================================
\section{Movimiento circular}

%------------------------------------------------
\begin{frame}{5. Movimiento circular}
Un punto que se mueve sobre una circunferencia se describe como:

\[
x=r\cos(\theta)
\]

\[
y=r\sin(\theta)
\]

donde:

\begin{itemize}
    \item \(r\) es el radio.
    \item \(\theta\) es el ángulo.
    \item \((x,y)\) es la posición del punto.
\end{itemize}
\end{frame}

%------------------------------------------------
\begin{frame}{Geometría del movimiento circular}
\begin{center}
\begin{tikzpicture}[scale=0.9]
    \draw[->] (-0.5,0) -- (4.8,0) node[right] {$x$};
    \draw[->] (0,-0.5) -- (0,3.5) node[above] {$y$};

    \coordinate (O) at (0,0);
    \coordinate (P) at (3.6,2.2);
    \coordinate (X) at (3.6,0);

    \draw[dashed] (O) circle (4.22);
    \draw[thick,blue] (O) -- (P) node[midway,above left] {$r$};
    \draw[thick,orange] (O) -- (X) node[midway,below] {$r\cos\theta$};
    \draw[thick,green!50!black] (X) -- (P) node[midway,right] {$r\sin\theta$};

    \draw[->,red,thick] (0.9,0) arc (0:31:0.9);
    \node[red] at (1.2,0.35) {$\theta$};

    \filldraw[red] (P) circle (2pt);
\end{tikzpicture}
\end{center}

Seno y coseno son proyecciones del radio sobre los ejes.
\end{frame}

%------------------------------------------------
\begin{frame}{Del movimiento circular a la oscilación}
Si observamos solo la coordenada \(y\):

\[
y=r\sin(\theta)
\]

\medskip

Y si el ángulo depende del tiempo:

\[
\theta(t)=\omega t
\]

entonces:

\[
y(t)=r\sin(\omega t)
\]

\medskip

\begin{block}{Idea clave}
Una oscilación es la proyección de un movimiento circular.
\end{block}
\end{frame}

%================================================
\section{Rotación en 2D}

%------------------------------------------------
\begin{frame}{6. Rotación en 2D}
Rotar un punto significa cambiar su ángulo sin cambiar su distancia al origen.

\[
\theta \longrightarrow \theta+\varphi
\]

\medskip

La distancia \(r\) permanece constante.

\medskip

La nueva posición se obtiene con:

\[
x' = x\cos\varphi-y\sin\varphi
\]

\[
y' = x\sin\varphi+y\cos\varphi
\]
\end{frame}

%------------------------------------------------
\begin{frame}{Matriz de rotación}
La rotación puede escribirse como:

\[
\begin{bmatrix}
x'\\
y'
\end{bmatrix}
=
\begin{bmatrix}
\cos\varphi & -\sin\varphi\\
\sin\varphi & \cos\varphi
\end{bmatrix}
\begin{bmatrix}
x\\
y
\end{bmatrix}
\]

\medskip

Esta matriz toma un punto y produce su versión rotada.
\end{frame}

%------------------------------------------------
\begin{frame}{Propiedades de la rotación}
Una rotación en 2D:

\begin{itemize}
    \item Conserva distancias.
    \item Conserva ángulos.
    \item No deforma el objeto.
    \item Cambia la orientación.
\end{itemize}

\medskip

Por eso se considera una transformación rígida.
\end{frame}

%------------------------------------------------
\begin{frame}{Rotación en animación}
Si el ángulo depende del tiempo:

\[
\varphi(t)=\omega t
\]

entonces la rotación cambia cuadro a cuadro.

\medskip

Aplicación:

\begin{itemize}
    \item Girar un objeto.
    \item Mover una articulación.
    \item Rotar una rueda.
    \item Orientar un avatar.
\end{itemize}
\end{frame}

%================================================
\section{Vectores}

%------------------------------------------------
\begin{frame}{7. Vectores}
Un vector en 2D se escribe como:

\[
\vec{v}=(v_x,v_y)
\]

\medskip

Un vector puede representar:

\begin{itemize}
    \item Desplazamiento.
    \item Dirección.
    \item Velocidad.
    \item Fuerza.
\end{itemize}

\medskip

En animación, los vectores son esenciales para mover objetos.
\end{frame}

%------------------------------------------------
\begin{frame}{Vector como desplazamiento}
Si un punto es:

\[
P=(x,y)
\]

y un vector de desplazamiento es:

\[
\vec{d}=(d_x,d_y)
\]

entonces:

\[
P'=P+\vec{d}
\]

\[
P'=(x+d_x,\;y+d_y)
\]

\medskip

Esto representa una traslación.
\end{frame}

%------------------------------------------------
\begin{frame}{Magnitud de un vector}
La magnitud mide la longitud del vector:

\[
\|\vec{v}\|=\sqrt{v_x^2+v_y^2}
\]

\medskip

Si \(\vec{v}\) representa velocidad, entonces:

\[
\|\vec{v}\|
\]

indica la rapidez total.

\medskip

La dirección está dada por la orientación del vector.
\end{frame}

%------------------------------------------------
\begin{frame}{Vector normalizado}
Un vector normalizado tiene longitud \(1\):

\[
\hat{v}=\frac{\vec{v}}{\|\vec{v}\|}
\]

\medskip

Se usa para separar:

\begin{itemize}
    \item Dirección.
    \item Rapidez.
\end{itemize}

\medskip

Movimiento típico:

\[
P_{\text{nuevo}}=P_{\text{actual}}+\hat{v}\cdot s\cdot dt
\]

donde \(s\) es la rapidez.
\end{frame}

%------------------------------------------------
\begin{frame}{Conclusión general}
Los conceptos estudiados muestran que animar es modelar matemáticamente el cambio.

\medskip

\[
\text{Animación}=\text{Funciones del tiempo}+\text{Geometría}
\]

\medskip

\begin{itemize}
    \item Las funciones del tiempo producen cambio.
    \item La velocidad produce movimiento.
    \item La integración acumula desplazamientos.
    \item El seno produce oscilación.
    \item Seno y coseno producen movimiento circular.
    \item Las matrices producen rotación.
    \item Los vectores organizan dirección y magnitud.
\end{itemize}
\end{frame}

\end{document}
